% ============================================================
% Sección: Impacto cuantificado y mitigación
% Responsable: Zambrano Triviño Leonel Leonidas
% Requiere en el preámbulo: \usepackage{booktabs,array}
% Citas: claves definidas en leonel.bib (pegar en el .bib del grupo)
% ============================================================

\section{Impacto cuantificado y mitigación}

\subsection*{Impacto cuantificado}

\begin{table}[h]
\centering
\small
\begin{tabular}{@{}>{\raggedright\arraybackslash}p{0.27\linewidth}>{\raggedright\arraybackslash}p{0.55\linewidth}l@{}}
\toprule
\textbf{Indicador} & \textbf{Dato} & \textbf{Fuente} \\
\midrule
Red afectada & 85\,\% de la red devolvió errores & \cite{fastly2021summary} \\
Clientes afectados & ``Casi todos'' los clientes de Fastly & \cite{fastly2021q2} \\
Detección & Menos de 1 minuto & \cite{fastly2021q2} \\
Recuperación parcial & 95\,\% de la red normal en 49 min & \cite{fastly2021summary} \\
Sitios caídos (ejemplos) & Amazon, Reddit, NYTimes, eBay, gov.uk, Twitch & \cite{medina2021thousandeyes,cnbc2021q2} \\
SLA & Créditos (compensaciones) a clientes; monto no publicado & \cite{fastly2021q2} \\
Ingresos Q2 2021 & 85,1 M USD frente a 85,73 M esperados & \cite{cnbc2021q2} \\
Pronóstico 2021 & Rebajado a 340--350 M USD por el incidente & \cite{fastly2021q2} \\
Clientes perdidos & Un cliente del top 10 no devolvió su tráfico & \cite{fastly2021q2} \\
\bottomrule
\end{tabular}
\caption{Impacto del incidente de Fastly (8 de junio de 2021).}
\end{table}

El costo directo no fue la caída de 49 minutos, sino lo que vino después: Fastly
cobra por uso, así que la baja de tráfico y los créditos por incumplimiento del
SLA (acuerdo de nivel de servicio) redujeron sus ingresos del segundo trimestre,
y la empresa bajó su pronóstico anual \cite{fastly2021q2}. Fastly no publicó una
cifra total de pérdidas; el equipo no usa estimaciones de terceros sin respaldo.

El impacto tampoco fue igual para todos los clientes. Reddit, que usaba solo
Fastly, estuvo caído toda la interrupción; NYTimes cambió su DNS hacia sus
servidores en Google Cloud y se recuperó en unos 40 minutos; Amazon, que reparte
su tráfico entre tres CDN, sacó a Fastly del camino hacia las 10:16 UTC
\cite{medina2021thousandeyes}.

\subsection*{Mitigación aplicada por la empresa}

\textbf{Medidas inmediatas.} El monitoreo detectó el fallo en menos de un minuto
y se publicó un aviso de estado a las 09:58 UTC \cite{fastly2021summary,fastly2021q2}.
A las 10:27 UTC ingeniería identificó la configuración del cliente que activaba el
defecto \cite{fastly2021summary} y la deshabilitó, con lo que la red empezó a volver
a la normalidad \cite{fastly2021tweet}. Con la red estable, se creó la corrección
definitiva y su despliegue comenzó a las 17:25 UTC \cite{fastly2021summary}.

\textbf{Cambios permanentes.} Fastly anunció cuatro acciones: desplegar la
corrección en toda la red de forma rápida y segura; hacer un \textit{post-mortem}
completo de los procesos seguidos; averiguar por qué el defecto no se detectó en
las pruebas y el control de calidad del despliegue del 12 de mayo; y reducir su
tiempo de remediación \cite{fastly2021summary}. En agosto reiteró que seguía
implementando medidas para aumentar la resiliencia \cite{fastly2021q2}.

\subsection*{Lecciones para el PFC de Leonel Zambrano}

% >>> AJUSTAR: reemplazar los nombres genéricos por los componentes reales del PFC
\begin{table}[h]
\centering
\small
\begin{tabular}{@{}>{\raggedright\arraybackslash}p{0.28\linewidth}>{\raggedright\arraybackslash}p{0.31\linewidth}>{\raggedright\arraybackslash}p{0.35\linewidth}@{}}
\toprule
\textbf{Riesgo} & \textbf{Consecuencia} & \textbf{Medida preventiva (TA-PFC-E1)} \\
\midrule
Desplegar una nueva versión en todos los nodos a la vez &
Un error latente tumba todos los nodos al mismo tiempo, como el defecto del 12 de mayo &
Despliegue escalonado: primero un nodo, verificar, luego el resto; con reversión (\textit{rollback}) preparada \\

Aceptar cambios de configuración sin validarlos &
Una configuración válida pero no probada activa un error y se propaga a todo el sistema &
Validar la configuración contra un esquema y probarla en un entorno de pruebas antes de aplicarla \\

Depender de un solo proveedor externo (nube, CDN o servicio de mensajería) &
Si el proveedor cae, cae todo el proyecto, como Reddit con Fastly &
Plan de respaldo documentado: proveedor o ruta alternativa y cambio de DNS, como hizo NYTimes \\
\bottomrule
\end{tabular}
\caption{Riesgos del PFC inspirados en el incidente de Fastly.}
\end{table}
